\section{Where the backdoor lives in a ViT}
\label{sec:mechanism}

\Cref{sec:results} is a ranking. This section explains it by following a backdoor through a ViT with 3 tools. It then checks that the picture they give predicts 3 things: the ranking, the reversal between noise and masking and the attacks on which the original paper's explanation holds. Every experiment runs on the same backdoored models as the detection results, with the benign model of the same dataset, probed with the same trigger, as the control.

\subsection{The backdoor direction and the depth at which it appears}
\label{sec:mechanism:direction}

For a model and a layer, take the mean difference between the class token's activation on a triggered image and on its clean copy. This is the backdoor direction of that layer, the construction Karayal\c{c}{\i}n et al.~\cite{karayalcin2026backdoor} use on ViTs. \Cref{tab:direction-ablation} tests whether it is causal. Projecting the final layer's direction out of the class token, after the final LayerNorm where the classifier reads, takes attack success to \AblationDirectionMinAsr{} on BadNets and leaves it at most \AblationDirectionMaxAsr{} on any of the \AblationCheckpoints{} attacks, while clean accuracy falls by at most \AblationCleanAccuracyCost{}.

2 controls say the effect is specific. Projecting out a random direction of the same norm leaves attack success at \AblationRandomDirectionMinAsr{} or above, and zeroing the 20 coordinates the trigger moves most leaves it at \AblationTopCoordinatesMinAsr{} or above. The backdoor is a direction in the residual stream rather than a set of neurons, and any method that ranks coordinates, heads or channels sees only a rotated shadow of it. LF is the exception at the final layer, where removing its direction takes attack success only to \AblationDirectionMaxAsr{}. Only when the direction is removed from every block's write to the stream (\cref{app:picek}) does LF fall to \ErasureLfAsrAfterAllWritesMax{}, so the low-frequency trigger is carried by the stream over several blocks rather than read off 1 layer.

% generated by scripts/paper/mech_direction.py from results/backdoor_neuron_ablation.json at 2026-09-23T11:19:21.201473+00:00, commit fe1619cdcc6be48ee0458803c6a7dd41e209c3bc-dirty. Do not edit by hand.
\begin{table}[htbp]
\centering
\small
\caption{The causal ablation at the peak layer, CIFAR-10 at the highest panel rate: attack success over clean accuracy after projecting out the rank-1 backdoor direction, a random direction, the top-20 trigger-activated coordinates and 20 random coordinates, as a forward hook after the final LayerNorm.}
\label{tab:direction-ablation}
\begin{adjustbox}{max width=\linewidth}
\begin{tabular}{lrrrrrrr}
\toprule
attack & peak layer & ASR & CA & rank-1 direction & random direction & top-20 coordinates & random-20 coordinates \\
\midrule
\bottomrule
\end{tabular}
\end{adjustbox}
\end{table}


The direction is not present from the start. \Cref{fig:tac-layers} follows its norm, relative to the clean activation's norm, through the 12 blocks on GTSRB, CIFAR-100 and Tiny ImageNet. On the benign model it never exceeds \DirectionBenignPeakRelNorm{}. On a backdoored model it stays near the benign level through the first half of the network, reaches half of its final value between blocks 7 and 9 depending on the attack, and peaks at blocks 9 to 12. The onset comes earlier for the global triggers than for the patch trigger on CIFAR-100 and Tiny ImageNet, and the order reverses on GTSRB, so the depth at which a trigger is written depends on the dataset as well as the trigger.

\begin{figure*}[t]
\centering
\includegraphics[width=\textwidth]{mech_tac_layers.pdf}
\caption{Relative norm of the backdoor direction (top) and mean trigger-activated change (TAC) \cite{zheng2022clp} per dimension, the mean absolute activation difference between a triggered image and its clean copy (bottom), by block, on GTSRB, CIFAR-100 and Tiny ImageNet at 10\% poisoning, 500 paired images per model. The benign model probed with the same trigger is dashed. The backdoor is written in the second half of the network.}
\label{fig:tac-layers}
\end{figure*}

\subsection{Routing of the trigger through attention}
\label{sec:mechanism:routing}

Activation patching \cite{zhang2024patching} asks where the backdoor sits at each depth. At a chosen block we overwrite a group of tokens on the triggered image with their values on the clean copy, then measure how much of the clean prediction comes back. \Cref{fig:patching} does this for the trigger's own tokens (the 4 patches under a BadNets trigger), for the class token and for a random token group of the same size.

For the local triggers, patching the trigger's own tokens recovers \PatchingLocalEarlyTrigger{} of the clean prediction over blocks \PatchingEarlyLayers{} and still recovers it at block \PatchingBadnetLastFull{}. The class token overtakes the trigger tokens only at block \PatchingLocalCrossover{} on average, and on BadNets it recovers \PatchingBadnetFinalCls{} at the last block. The random group never exceeds \PatchingRandomNullMax{} at any depth, which is the null this rests on. For a local trigger the backdoor therefore sits entirely in the trigger's own tokens for most of the network and reaches the class token, where the classifier reads, only in the last quarter. A global trigger covers every patch, so its trigger set is all \VitPatches{} patch tokens and the random group of the same size is the same set. For these attacks the experiment can only confirm that the backdoor is spread over the tokens.

\begin{figure*}[t]
\centering
\includegraphics[width=\textwidth]{mech_activation_patching.pdf}
\caption{Share of the clean prediction recovered when a token group is replaced by its clean value at a given block, averaged over the backdoored models of each attack. For the local triggers (BadNets, TaCT) the trigger's tokens hold the backdoor through block 10 and the class token takes over only at the end.}
\label{fig:patching}
\end{figure*}

The class token's attention confirms the route. On the BadNets models the share of its attention that falls on the trigger's 4 tokens is \RoutingBadnetClsEarly{} in blocks 1 to 4, \RoutingBadnetClsRouting{} in blocks 5 to 8 and \RoutingBadnetClsLate{} in blocks 9 to 12, against \RoutingBenignClsLate{} on the benign model probed with the same trigger, over \RoutingCells{} checkpoints. Attention, the only operation that moves information between tokens, carries the trigger to the class token in the second half of the network. \Cref{fig:routing} shows this per block.

\begin{figure*}[t]
\centering
\includegraphics[width=0.9\textwidth]{mech_routing.pdf}
\caption{Left, share of the class token's attention that falls on the trigger's tokens, by block and attack, against the benign model probed with the same trigger (dashed). Right, the norm rank of the trigger's tokens among all patch tokens. The trigger is routed to the class token in blocks 5 to 12.}
\label{fig:routing}
\end{figure*}

\subsection{What this means for the placement}
\label{sec:mechanism:placements}

Put together, the 3 tools give a simple account. A backdoored input's decision rests on a few tokens whose content is moved to the class token late through attention. A clean input's decision rests on many tokens. Masking whole tokens at the attention input, before the mixing happens, removes the trigger's tokens with a probability that grows with the mask rate. When it does, the shortcut is gone and the prediction moves. On a clean input the remaining tokens still carry the class and the prediction stays.

This is why PSBD-TM separates the 2 populations, and why it should act in every block rather than a band. The mask must be in place at whatever depth the routing happens. The same mask at the MLP input acts after the block's attention has already mixed the trigger into the class token, so it removes less of the shortcut, which is the \StaircaseOperatorsBeforeMlpNormTokenMaskGain{} cost of moving the site in \cref{tab:staircase-operators}. Dropout on the residual stream after the add perturbs the class token itself, on clean and triggered inputs alike, which is why PSBD-RD reads a coin on the attacks whose shortcut is weak.

\subsection{LayerNorm and additive noise}
\label{sec:mechanism:layernorm}

The reversal in \cref{sec:results:operator} has a mechanical cause. Every perturbation at the attention input passes through a LayerNorm before it reaches the attention. A LayerNorm rescales each token to unit variance, so an additive perturbation that raises the token's variance is partly divided away. \Cref{tab:layernorm-absorption} measures it. The attention input's LayerNorm passes \AbsorptionAttentionInputGaussianSurvival{} of a Gaussian disturbance and \AbsorptionAttentionInputTokenMaskSurvival{} of a token mask, which sets whole tokens to 0 and moves a token's variance very little. At the positions no LayerNorm follows every operator survives at \AbsorptionControlSurvival{}, which is the control that this is the norm's doing. The measurement is on \AbsorptionModels{} models, so it establishes the mechanism rather than its spread.

% generated by scripts/paper/mech_layernorm.py from results/*/layernorm_absorption.json (2 models) at 2026-09-23T16:51:28.045898+00:00, commit d4a8014354b8fcc8d014beec50fef60ed46fa47d-dirty. Do not edit by hand.
\begin{table}[htbp]
\centering
\small
\caption{What a LayerNorm passes through. Survival is the perturbation's relative norm after the LayerNorm over its relative norm before it, averaged over blocks and over 2 models. The last 2 rows sit after the LayerNorm, where nothing normalizes them, and are the control.}
\label{tab:layernorm-absorption}
\begin{adjustbox}{max width=\linewidth}
\begin{tabular}{llrrr}
\toprule
position & a norm follows & token mask & channel mask & gaussian \\
\midrule
attention input & yes & 1.002 & 0.968 & 0.821 \\
MLP input & yes & 0.971 & 0.981 & 0.872 \\
attention input after norm & no & 1.000 & 1.000 & 1.000 \\
MLP input after norm & no & 1.000 & 1.000 & 1.000 \\
\bottomrule
\end{tabular}
\end{adjustbox}
\end{table}


Noise at the attention input is therefore weaker than its nominal rate says, and loses \StaircaseOperatorsBeforeAttentionNormGaussianGain{} to masking. At the MLP input the same noise is applied per token after the attention has done its mixing. There masking a token no longer removes routed evidence, and noise injected after that sublayer's LayerNorm reads \StaircaseOperatorsBeforeMlpGaussianAuroc{} against \StaircaseOperatorsBeforeMlpNormTokenMaskAuroc{} for a mask injected before it. The reversal is a consequence of what sits behind each site.

\subsection{The prediction shift phenomenon on ViT}
\label{sec:mechanism:phenomenon}

The PSBD paper explains its statistic by neuron bias. Under dropout the network falls back on the classes its neurons are biased toward, and because the backdoor's neurons are biased toward the target most firmly, clean inputs whose prediction shifts should land on the target class. We measured this directly on every model. Among clean validation images whose prediction changes under PSBD-TM, what share lands on the attacker's target class? \Cref{fig:shift-target} plots that share against the achieved shift ratio, per attack, with the benign models' band as the control, and \cref{tab:mech-shift-target-by-dataset} gives it per dataset.

\begin{figure}[t]
\centering
\includegraphics[width=\linewidth]{mech_shift_target.pdf}
\caption{Share of shifted clean predictions that land on the target class, against the achieved clean shift ratio, 1 curve per attack averaged over models. The gray band is the benign range and the dashed line the uniform expectation.}
\label{fig:shift-target}
\end{figure}

% generated by scripts/paper/mech_shift_target.py from results/<folder>/psbd_metrics.json (65 clearing cells, 4 benign) at 2026-09-23T16:54:56.873920+00:00, commit d4a8014354b8fcc8d014beec50fef60ed46fa47d-dirty. Do not edit by hand.
\begin{table}[htbp]
\centering
\small
\caption{Share of shifted clean predictions that land on the target class under PSBD-TM, per dataset and attack, against the uniform expectation and the benign reference of the dataset.}
\label{tab:mech-shift-target-by-dataset}
\begin{adjustbox}{max width=\linewidth}
\begin{tabular}{lrrrrrr}
\toprule
Dataset & Attack & n & Share & Uniform & Benign & Excess \\
\midrule
CIFAR-10 & BPP & 3 & 0.166 & 0.100 & 0.004 & +0.162 \\
CIFAR-10 & WaNet & 2 & 0.172 & 0.100 & 0.004 & +0.168 \\
CIFAR-10 & TaCT & 3 & 0.000 & 0.100 & 0.004 & $-$0.004 \\
CIFAR-10 & SIG & 1 & 0.000 & 0.100 & 0.004 & $-$0.004 \\
CIFAR-10 & BadNets & 3 & 0.518 & 0.100 & 0.004 & +0.514 \\
CIFAR-10 & Blend & 3 & 0.005 & 0.100 & 0.004 & +0.001 \\
CIFAR-10 & LF & 3 & 0.172 & 0.100 & 0.004 & +0.168 \\
CIFAR-100 & BPP & 3 & 0.157 & 0.010 & 0.000 & +0.157 \\
CIFAR-100 & TaCT & 3 & 0.004 & 0.010 & 0.000 & +0.004 \\
CIFAR-100 & BadNets & 3 & 0.002 & 0.010 & 0.000 & +0.002 \\
CIFAR-100 & Blend & 3 & 0.009 & 0.010 & 0.000 & +0.009 \\
CIFAR-100 & LF & 3 & 0.256 & 0.010 & 0.000 & +0.256 \\
GTSRB & BPP & 3 & 0.001 & 0.023 & 0.000 & +0.001 \\
GTSRB & WaNet & 1 & 0.996 & 0.023 & 0.000 & +0.996 \\
GTSRB & TaCT & 2 & 0.000 & 0.023 & 0.000 & +0.000 \\
GTSRB & BadNets & 3 & 0.793 & 0.023 & 0.000 & +0.793 \\
GTSRB & Blend & 3 & 0.015 & 0.023 & 0.000 & +0.015 \\
GTSRB & LF & 3 & 0.054 & 0.023 & 0.000 & +0.054 \\
Tiny ImageNet & BPP & 3 & 0.347 & 0.005 & 0.002 & +0.346 \\
Tiny ImageNet & WaNet & 2 & 0.400 & 0.005 & 0.002 & +0.398 \\
Tiny ImageNet & TaCT & 3 & 0.002 & 0.005 & 0.002 & +0.000 \\
Tiny ImageNet & BadNets & 3 & 0.150 & 0.005 & 0.002 & +0.149 \\
Tiny ImageNet & Blend & 3 & 0.005 & 0.005 & 0.002 & +0.003 \\
Tiny ImageNet & LF & 3 & 0.385 & 0.005 & 0.002 & +0.383 \\
\bottomrule
\end{tabular}
\end{adjustbox}
\end{table}


The phenomenon is real and it is selective. On GTSRB, with 43 classes and a local trigger, 69\% (BadNets) and 93\% (WaNet) of shifted clean predictions land on the target, far above the uniform 2.3\% and the benign model's 0\%. On CIFAR-10 BadNets it is 36\%. For Blend and TaCT it is at or below the uniform rate on every dataset, and for BadNets on CIFAR-100 it is 0.3\% against a uniform 1\%. Over all models the mean excess over benign is +0.118 with an interval from +0.066 to +0.179, and it is carried entirely by the local triggers on the few-class datasets.

Detection works in both regimes. Blend is among the best-detected attacks on every dataset with an excess of 0. The neuron-bias account therefore describes what happens when a trigger is routed through few tokens toward a target that is easy to fall into, and the margin appears to be what makes the statistic work in general. A triggered input's predicted class has a large margin over the rest because the shortcut is firm, and any perturbation that moves clean decisions more than firm ones separates the populations, whether or not the clean decisions move toward the target.

If that were the whole story, the unperturbed confidence would separate the populations by itself and the perturbation would be doing no work. It does not. On \ConfidenceNullCheckpoints{} backdoored models the maximum softmax probability reads a mean AUROC of \ConfidenceNullAuroc{} against \ConfidenceNullPsuAuroc{} for PSU under PSBD-TM, and PSU is ahead on \ConfidenceNullPsuWins{} of them. On the benign references PSU reads \ConfidenceNullBenignPsu{}, which is the check that it is measuring the backdoor rather than the architecture. What the statistic reads is how robust a firm decision is, not how confident it started.

\subsection{Agreement with Karayal\c{c}{\i}n et al.}
\label{sec:mechanism:picek}

Karayal\c{c}{\i}n et al.~\cite{karayalcin2026backdoor} studied backdoor directions in ViT-B/16 on the same 3 datasets with an overlapping attack set, and our models let us test 3 of their claims. Their removal of the backdoor by projecting the direction out of every weight matrix that writes to the residual stream reproduces on our models. Over \ErasureCheckpoints{} checkpoints it takes attack success from \ErasureAsrBeforeMean{} to \ErasureAsrAfterAllWritesMean{} at a clean-accuracy cost of \ErasureCaBeforeMean{} to \ErasureCaAfterAllWritesMean{}. Editing the last 2 blocks alone leaves \ErasureAsrAfterBlocksOneZeroOneOneMean{} and removing a random direction leaves \ErasureAsrAfterRandomDirectionMean{}, so neither control moves the attack. Blend on CIFAR-100 survives it exactly as they report, which is the \ErasureAsrAfterAllWritesMax{} in that column. The direction is re-created by every block, so removing it means removing it everywhere, which is also why the perturbation that acts in every block wins in \cref{sec:results:depth}. Their observation that distributed triggers reach the class token earlier than patch triggers holds on CIFAR-100 and Tiny ImageNet and reverses on GTSRB. Their data-free weight detector reproduces on WaNet and only in part on BPP, for the reason \cref{tab:direction-ablation} gives, since a score built from coordinates of the head and the weights sees only the projection of a rotated direction. Their projected gradient descent readings reproduce in part, with the high mid-layer cosine between an adversarial shift and the backdoor direction holding for BPP, Blend and LF but not WaNet. \Cref{app:picek} gives the full numbers.
